% Generated draft: requires mathematical and source review.
\documentclass[11pt]{article}
\usepackage[margin=25mm]{geometry}
\usepackage{fontspec}
\setmainfont{Latin Modern Roman}
\usepackage{amsmath,amssymb,mathtools,unicode-math}
\setmathfont{Latin Modern Math}
\usepackage{enumitem,needspace,xurl,hyperref}
\setlength{\parindent}{0pt}
\setlength{\parskip}{5pt}
\setlength{\emergencystretch}{3em}
\newcommand{\N}{\mathbb{N}}
\newcommand{\Z}{\mathbb{Z}}
\newcommand{\Q}{\mathbb{Q}}
\newcommand{\R}{\mathbb{R}}
\newcommand{\C}{\mathbb{C}}
\newcommand{\F}{\mathbb{F}}
\newcommand{\A}{\mathbb{A}}
\newcommand{\PP}{\mathbb{P}}
\newcommand{\E}{\mathbb{E}}
\newcommand{\eps}{\varepsilon}
\newcommand{\dd}{\,\mathrm{d}}
\begin{document}
\section*{UnsolvedMath ID: 30000686}
\textbf{Source problem code:} catalogue entry
\subsection{Definitions and mathematical statement}
There is a substantive conflict between the original and cleaned source formulations. Let \((W,S)\) be a Coxeter system, with Bruhat order \(\leq\), and let \(\theta\) be an involutive automorphism. The two subsets appearing in the record are
\[
\iota(\theta)=\{\theta(w^{-1})w:w\in W\}
\qquad\text{and}\qquad
\mathcal I(\theta)=\{x\in W:\theta(x)=x^{-1}\}.
\]
The first is the twisted-identity set in the original extraction, while the second is the twisted-involution set substituted in the cleaned text. They are not identical in general: when \(\theta=\operatorname{id}\), the first contains only the identity and the second contains all involutions. Therefore the conflict is not a mere typography correction.

For either chosen set \(P\), let \(\operatorname{Br}(P)\) denote its induced Bruhat-order subposet. For \(u<v\) in it, the usual interval-topology object is the order complex of the open interval
\[
(u,v)_P=\{x\in P:u<x<v\},
\]
whose simplices are finite chains. Such intervals are finite because ambient Bruhat intervals are finite. If the endpoints are included, their presence makes the order complex a cone, so the open-interval or proper-part convention is essential to the intended topological question. The source does not specify whether homeomorphism type, homotopy type, homology, or another invariant is sought.

The original passage asks two tasks for \(\operatorname{Br}(\iota(\theta))\): classify the Coxeter data for which it is graded, and describe its interval topology. Graded means that a rank function increases by one across every cover, equivalently the corresponding finite intervals have consistent maximal-chain lengths. The cleaned statement instead asks only for the interval topology of \(\operatorname{Br}(\mathcal I(\theta))\), and its supplied verification says that the report's Problem 5 concerns twisted involutions. Both formulations are retained as conflicting source evidence; the original grading question is not discarded and no theorem about one poset is used to answer the other.

The record also does not explicitly say \(\theta(S)=S\), the usual Coxeter-system automorphism hypothesis relevant to Bruhat-order structure. An arbitrary group automorphism and a diagram automorphism should not be identified silently. The supplied literature triage checked 2026-08-21 gives only an aggregate open assessment, which cannot resolve the discrepancy or certify a status for each distinct poset problem.
\subsection{Short English statement}
The original asks for gradedness and interval topology of twisted identities, while the cleaned version asks only about twisted involutions. These are genuinely different sets, so both formulations and the missing Coxeter-automorphism and interval conventions require source reconciliation.
\subsection{Sources}
\begin{itemize}
\item Local UnsolvedMath registry, permanent ID 30000686, catalogue entry, \emph{Topology of Twisted-Involution Bruhat Intervals}; its original and cleaned statements give different underlying sets.
\item \emph{Geometric and Topological Combinatorics} (2007), Oberwolfach Reports, pp. 195--272, DOI \url{https://doi.org/10.4171/owr/2007/04}.
\item Supplied verification identifies Problem 5 as an interval-topology question for twisted involutions, in conflict with the reproduced original twisted-identity and gradedness passage.
\item Supplied literature triage checked 2026-08-21 cites only that report; no separate theorem or reconciliation is supplied.
\end{itemize}
\end{document}
